\documentclass{article}
\usepackage{mathtools} 
\usepackage{fontspec}
\usepackage[UTF8]{ctex}
\usepackage{amsthm}
\usepackage{mdframed}
\usepackage{xcolor}
\usepackage{amssymb}
\usepackage{amsmath}


% 定义新的带灰色背景的说明环境 zremark
\newmdtheoremenv[
  backgroundcolor=gray!10,
  % 边框与背景一致，边框线会消失
  linecolor=gray!10
]{zremark}{说明}

% 通用矩阵命令: \flexmatrix{矩阵名}{元素符号}{行数}{列数}
\newcommand{\flexmatrix}[4]{
  \[
  #1 = \begin{pmatrix}
    #2_{11}     & #2_{12}     & \cdots & #2_{1#4}   \\
    #2_{21}     & #2_{22}     & \cdots & #2_{2#4}   \\
    \vdots      & \vdots      & \ddots & \vdots     \\
    #2_{#31}    & #2_{#32}    & \cdots & #2_{#3#4}
  \end{pmatrix}
  \]
}

% 简化版命令（默认矩阵名为A，元素符号为a）: \quickmatrix{行数}{列数}
\newcommand{\quickmatrix}[2]{\flexmatrix{A}{a}{#1}{#2}}

\begin{document}
\title{习题 1.2}
\author{张志聪}
\maketitle

\section*{7}

\begin{itemize}
  \item (1)

        (i)对任意$y \in A + B$，可以表示成
        \begin{align*}
          y = x_1 + x_2, x_1 \in A, x_2 \in B
        \end{align*}
        因为
        \begin{align*}
          x_1 \leq \sup A \\
          x_2 \leq \sup B
        \end{align*}
        所以，
        \begin{align*}
          y = x_1 + x_2 \leq \sup A + \sup B
        \end{align*}
        可得，$\sup A + \sup B$是$A + B$的一个上界。

        (ii)对$\forall \epsilon > 0$，存在$x_0 \in A, y_0 \in B$，使得
        \begin{align*}
          x_0 \geq \sup A - \frac{1}{2}\epsilon \\
          y_0 \geq \sup B - \frac{1}{2}\epsilon
        \end{align*}
        所以
        \begin{align*}
          x_0 + y_0 \geq \sup A + \sup B - \epsilon
        \end{align*}
        又因为
        \begin{align*}
          x_0 + y_0 \in A + B
        \end{align*}

        综上，$\sup A + \sup B$是$A + B$的最小上界，即
        \begin{align*}
          \sup (A + B) = \sup A + \sup B
        \end{align*}

\end{itemize}


\end{document}